\section{The dual of pathwidth}
\label{app:dual}

Every lower bound this paper uses reads one object at a time: a clique, a
vertex of large minimum degree in a minor, a prefix of a layout in Harper's
bound, a contracted copy of the graph in Theorem~\ref{thm:minor}, a minor
found by the search of Appendix~\ref{app:tabu}. Each is a certificate that
the pathwidth is at least some value, and each family has a ceiling on what
it can certify. This appendix asks the question behind all of them: is there
a single family with no ceiling, a \emph{dual} of pathwidth in the sense that
linear programming has a dual, and if so, why can it not simply be maximised?
The answer to the first question is yes, the blockages of
\citet{BRST1991}; the answer to the second is three theorems and one
conjecture. Throughout, $n$ is the order of a blockage, never the number of
vertices, which is written $|V|$. The definitions and the weak half of the
duality are formalised in \texttt{Blockage.lean} and
\texttt{BlockageClass.lean}, with no
\texttt{sorry}~\leanref{blockage}~\leanref{blockageclass}.

\subsection{The problem of the dual}

Pathwidth is a minimum: over all orderings of the vertices, the largest
number of vertices that a prefix has in common with the neighbourhood of its
complement, minimised. An ordering is therefore a short proof of an upper
bound: $|V|$ numbers, checked in linear time, and the search of
Section~\ref{sec:search} finds one in milliseconds on most instances. A lower
bound has no such object. The statement $\pw(G)\ge w$ is proved, in the body
of this paper, by exhausting a search at width $w-1$, and the proof is the
search tree, $10^{11}$ nodes on the hardest instances of Section~\ref{sec:practice}.

Linear programming resolves the same asymmetry by a second optimisation
problem: a maximisation whose every feasible point is a lower bound on the
minimum (weak duality) and whose optimum equals it (strong duality). A
feasible dual point is then a certificate of the lower bound, and the best
certificate is tight. For treewidth the dual object is the bramble of
\citet{SeymourThomas1993}: a family of connected, pairwise touching vertex
sets, whose order is the size of the smallest set meeting all of them, and
treewidth plus one is the largest order of a bramble. For pathwidth it is the
blockage. This appendix states the object, proves that it is the dual, and
then proves why maximising it is not an algorithm: not in polynomial time
unless $\mathrm P=\mathrm{NP}$, not with a short certificate unless
$\mathrm{NP}=\mathrm{coNP}$, and not by the two concrete dual searches one
would write first, whose cost is provably at least a binomial where the
primal's is small.

\subsection{Blockages}

Let $G=(V,E)$ be a finite simple graph. For $X\subseteq V$ the
\emph{attachments} of $X$ are its vertices with a neighbour outside $X$,
\[
\att(X) \;=\; \{\,v\in X : \exists u\notin X,\ uv\in E\,\},
\]
and $\alpha(X)=|\att(X)|$. Two sets $X_1,X_2$ are \emph{complementary} if
$X_1\cup X_2=V$ and $\att(X_1)\subseteq X_2$; equivalently no edge joins
$X_1\setminus X_2$ to $X_2\setminus X_1$, so a complementary pair is a
separation of $G$ with separator $X_1\cap X_2$, and every $X$ has the
canonical partner $X^{c}=(V\setminus X)\cup\att(X)$.

\begin{definition}[blockage, \citealp{BRST1991}]\label{def:blockage}
A \emph{blockage of order $n$} is a family $\mathcal B$ of subsets of $V$
such that
\begin{enumerate}
\item[(i)] every $X\in\mathcal B$ has $\alpha(X)\le n$;
\item[(ii)] if $X\in\mathcal B$, $Y\subseteq X$ and $\alpha(Y)\le n$, then $Y\in\mathcal B$;
\item[(iii)] if $X_1,X_2$ are complementary and $|X_1\cap X_2|\le n$, then
  exactly one of $X_1,X_2$ lies in $\mathcal B$.
\end{enumerate}
\end{definition}

Read as a game, $\mathcal B$ is the strategy of an invisible fugitive against
$n+1$ searchers: for every way the searchers can cut the graph with at most
$n$ of them, (iii) says which side the fugitive is on, and (i), (ii) make
those choices consistent; this is Theorem~(5.1) of the paper. The bounds are
$\le n$ throughout, as in the paper's proof of its Theorem~(2.1) and in the
Lean statement; with strict bounds the family
$\{\emptyset,\{a\},\{b\},\{c\},\{a,b\}\}$ on the path $a$--$b$--$c$ would be a
``blockage of order 2'' on a graph of pathwidth~1.

A blockage is a large object, but it is determined by a small part of
itself. Call $X$ \emph{big} if $\alpha(X)\le n$ and $X\notin\mathcal B$.

\begin{lemma}[a blockage is an avoidance family]\label{lem:frame}
Let $\mathcal B$ be a blockage of order $n$. Bigness is upward closed among
the sets with $\alpha\le n$, and
\[
\mathcal B \;=\; \{\,X\subseteq V : \alpha(X)\le n\ \text{and no big set is contained in } X\,\}.
\]
Hence $\mathcal B$ is determined by its \emph{frame}, the $\subseteq$-minimal
big sets $F_1,\dots,F_r$. For such a frame, axioms (i) and (ii) hold
automatically, and (iii) reads:
\begin{enumerate}
\item[(a)] every separation of order at most $n$ has a side containing some $F_j$;
\item[(b)] no separation of order at most $n$ has $F_i$ inside one side and $F_j$ inside the other.
\end{enumerate}
\end{lemma}
\begin{proof}
If $X$ is big and $X\subseteq Y$ with $\alpha(Y)\le n$, then $Y\notin\mathcal B$,
for otherwise (ii) would put $X$ in $\mathcal B$; so bigness is upward
closed. A set with $\alpha\le n$ is in $\mathcal B$ or big, and it is big
exactly when it contains a minimal big set, which gives the display. For the
reading of (iii): a complementary pair $(X_1,X_2)$ with $|X_1\cap X_2|\le n$
has $\att(X_i)\subseteq X_1\cap X_2$, so both sides satisfy (i); ``exactly
one in $\mathcal B$'' says exactly one side contains no frame set, which is
(a) and (b) together. Condition (b) is a Menger condition and is checked
pair by pair in polynomial time; condition (a) is the whole difficulty.
\end{proof}

Three frames are proved to give blockages in \texttt{BlockageClass.lean}: a
clique $K$ of $n+1$ vertices, with the single frame set $K$
(\texttt{isBlockageIn\_clique}); the pullback of a blockage of $H$ along a
minor model of $H$ in $G$, each frame set pulled back along its branch sets
(\texttt{isBlockageIn\_minor}); and three disjoint connected branches each
carrying a frame of order $k$, pairwise joined by connected connectors
missing the third, whose union carries a frame of order $k+1$
(\texttt{isBlockageIn\_branch}), the blockage form of the corrected branch
lemma of Section~\ref{sec:minor}. Every statement was checked by exhaustive
enumeration on every graph with at most six vertices before it was
stated~\leanref{blockageclass}.

\subsection{Why it is the dual}

\begin{theorem}[weak duality: a blockage is a lower bound]\label{thm:blockage}
If $G$ has a blockage of order $n$, then $\pw(G)\ge n$~\leanref{blockage}.
\end{theorem}
\begin{proof}
Suppose $\mathcal B$ is a blockage of order $n$ and, for a contradiction,
that $W_0,\dots,W_r$ is a path decomposition of $G$ of width less than $n$,
so every bag has at most $n$ vertices. Let $X_i=W_0\cup\dots\cup W_{i-1}$
be the prefix unions, with $X_0=\emptyset$ and $X_{r+1}=V$. Two facts follow
from the interval property of a path decomposition: $\att(X_i)\subseteq W_i$,
since a vertex of $X_i$ with a neighbour outside $X_i$ shares a bag with that
neighbour, which lies at or after $W_i$, so the vertex's interval of bags
reaches $W_i$; and $X_{i+1}\setminus X_i\subseteq W_i$ by definition. Hence
$\alpha(X_i)\le n$ for every $i$.

Now $\emptyset\in\mathcal B$ and $V\notin\mathcal B$, because $(\emptyset,V)$
is a complementary pair with empty intersection, and (iii) puts exactly one of
the two in $\mathcal B$, while (ii) with $Y=\emptyset$ rules out the second
choice. Lemma~(2.2) of the paper, proved in Lean as
\texttt{IsBlockage.mem\_of\_mem}, says: if $X\in\mathcal B$, $\alpha(Y)\le n$ and
$|(Y\setminus X)\cup\att(X)|\le n$, then $Y\in\mathcal B$. Apply it to
$X=X_i\in\mathcal B$ and $Y=X_{i+1}$: $\alpha(X_{i+1})\le n$, and
$(X_{i+1}\setminus X_i)\cup\att(X_i)\subseteq W_i$ has at most $n$
vertices, so $X_{i+1}\in\mathcal B$. By induction $V=X_{r+1}\in\mathcal B$, a
contradiction.
\end{proof}

\begin{theorem}[strong duality, \citealp{BRST1991}, Theorem~(2.1)]\label{thm:blockage-strong}
$G$ has a blockage of order $n$ if and only if $\pw(G)\ge n$. Consequently
\[
\pw(G)\;=\;\max\{\,n : G \text{ has a blockage of order } n\,\},
\]
and the feasible orders form an interval: a blockage of order $n$ restricted
to its members with $\alpha\le n-1$ is a blockage of order $n-1$.
\end{theorem}

The ``if'' direction is the paper's Sections~3 and~4, through partial
blockages and fractures, and it is not formalised here. A shorter route,
which we record because it is the one a formalisation would take, goes
through orientations: \citet{Erde2017}, Theorem~12, proves by induction on
the set of separations that a graph has either a path decomposition of
small width or a consistent orientation of its small separations with no two
chosen sides covering the graph, never both; and \citet{DiestelOum2019},
Theorem~7.2, shows that such an orientation is exactly a blockage, the
``two sides'' clause being the paper's stoppage condition. Strengthening
``two'' to ``three'' gives the tangle axioms, which is why blockages are
stronger than tangles and $\pw\ge\tw$. The restriction claim is immediate
from the axioms: the separations of order at most $n-1$ are among those of
order at most $n$, and both sides of such a separation have $\alpha\le n-1$.

Weak duality is what makes a bottom-up search for lower bounds possible, and
strong duality is what says such a search has no ceiling. Every lower-bound
family in this paper is a restricted class of blockages. A clique of size
$n+1$ gives one, with the clique as its only frame set. Harper's
vertex-isoperimetric bound at one prefix size is a blockage built from the
separations at that size alone. A minor $H$ with $\pw(H)\ge n$ gives a
blockage of $G$ by pulling $H$'s back along the branch sets, which is
Theorem~\ref{thm:minor} read from the dual side. The duality gaps measured in
Section~\ref{sec:certs}, family by family, are the gaps of these restricted
classes, and the finding there that the residual's obstruction is a band of
consecutive prefix sizes says, in this vocabulary, that the blockages of those
instances need separations at many consecutive orders at once, which no
single-prefix class supplies.

\subsection{Three theorems on maximising the dual}

Strong duality invites an algorithm: maximise the order over blockages, from
below, and stop when the order meets the width of the best layout. The
following three results say, in increasing concreteness, why that is not an
algorithm. The first two are standard complexity theory applied to one
hardness theorem; we state them because they are exactly what the dual is
being asked to do.

\begin{theorem}[no efficient maximisation]\label{thm:dual-np}
If there is a polynomial-time algorithm that, given $G$, returns the largest
order of a blockage of $G$, then $\mathrm P=\mathrm{NP}$. The same holds for
an algorithm that returns, for given $G$ and $n$, whether a blockage of order
$n$ exists.
\end{theorem}
\begin{proof}
By Theorem~\ref{thm:blockage-strong} the largest order of a blockage is
$\pw(G)$, and a blockage of order $n$ exists if and only if $\pw(G)\ge n$. So
either algorithm decides, in polynomial time, whether $\pw(G)\le k$ for a
given $k$. That problem is NP-complete: \citet{LY2002ref5} prove it in the
form of interval graph completion, and \citet{LY2002ref14}, Theorem~7, in
the form of the vertex separation game, which equals pathwidth by
\citet{LY2002ref13}. A polynomial-time algorithm for an NP-complete problem
gives $\mathrm P=\mathrm{NP}$.
\end{proof}

\begin{theorem}[no short certificate]\label{thm:dual-conp}
Suppose there is a polynomial $p$ and a polynomial-time verifier $\mathcal V$
such that for every graph $G$ and every $n\le\pw(G)$ there is a string $c$ of
length at most $p(|V|)$ with $\mathcal V(G,n,c)=1$, while
$\mathcal V(G,n,c)=0$ for every $c$ whenever $n>\pw(G)$. Then
$\mathrm{NP}=\mathrm{coNP}$. In particular, no representation of blockages
can be both polynomial in size and polynomial to check on every graph unless
$\mathrm{NP}=\mathrm{coNP}$.
\end{theorem}
\begin{proof}
Let $L=\{(G,n):\pw(G)\ge n\}$. Its complement, $\{(G,k):\pw(G)\le k\}$ with
$k=n-1$, is NP-complete by the theorems cited above, so $L$ is
coNP-complete: every language in coNP reduces to $L$ in polynomial time. The
verifier $\mathcal V$ with its short strings $c$ is by definition a
nondeterministic polynomial-time acceptor of $L$, so $L\in\mathrm{NP}$. Then
every coNP language reduces to a language in NP, and NP is closed under
polynomial-time reductions, so $\mathrm{coNP}\subseteq\mathrm{NP}$; taking
complements, $\mathrm{NP}\subseteq\mathrm{coNP}$. The last sentence is the
case in which $c$ encodes a blockage of order $n$ in any fixed representation
and $\mathcal V$ checks the axioms
\citep[Chapter~7]{GareyJohnson1979}.
\end{proof}

The theorem is conditional and cannot be otherwise: an unconditional proof
that blockages have no short representation would separate NP from coNP.
What it rules out is representation-independent. A frame of $r$ sets
(Lemma~\ref{lem:frame}) is a certificate whose verifier runs in time
$|V|^{O(r)}$, polynomial for bounded $r$, and the theorem says that no
scheme makes $r$ bounded on every graph. For treewidth the same statement is
in print: \citet{GroheMarx2009} write that a bramble ``is not a good
characterization of tree width in the coNP sense'' for exactly this reason,
and prove the unconditional form in a fixed representation, to which we
return in Conjecture~\ref{conj:blockage-size}. We found no statement of
either kind for pathwidth and blockages.

The third theorem is about the two dual searches one would actually write,
and it is unconditional. Fix a width $w$ and let
\[
\mathcal F_w \;=\; \{\,S\subseteq V : |N(S)\setminus S|\le w\,\},
\]
the sets whose outer boundary has at most $w$ vertices. A layout of width at
most $w$ is a chain $\emptyset=S_0\subset S_1\subset\dots\subset S_{|V|}=V$
of sets in $\mathcal F_w$, each obtained from the last by adding one vertex,
so $\pw(G)\le w$ exactly when $\emptyset$ and $V$ lie in the same connected
component of the graph on $\mathcal F_w$ whose edges are single additions.
Our corrected algorithm explores the component of $\emptyset$; its states
are the sets of that component, within a factor of about two
(Section~\ref{sec:certs}). The natural dual search explores the component of
$V$ instead, removing vertices one at a time; a refutation by either search
is the exhaustion of its component.

\begin{theorem}[the lattice floor of a top-down search]\label{thm:dual-lattice}
Let $w<|V|$. Every set $S$ with $|V\setminus S|\le w$ lies in $\mathcal F_w$
and in the component of $V$. Hence a refutation of width $w$ that exhausts
the component of $V$ visits at least $\sum_{j=0}^{w}\binom{|V|}{j}$ states,
on every graph. Moreover, complementation $S\mapsto V\setminus S$ is an
order-reversing bijection from $\mathcal F_w$ onto
$\mathcal F'_w=\{T : \alpha(T)\le w\}$, the sets whose inner boundary has at
most $w$ vertices, carrying single additions to single removals; so the
top-down search on $\mathcal F'_w$ is the bottom-up search on
$\mathcal F_w$, and the two directions that avoid the binomial region are
the same algorithm.
\end{theorem}
\begin{proof}
If $|V\setminus S|\le w$ then $N(S)\setminus S\subseteq V\setminus S$ has at
most $w$ vertices, so $S\in\mathcal F_w$; and removing the vertices of
$V\setminus S$ from $V$ one at a time passes only through sets with still
smaller complements, each in $\mathcal F_w$ by the same argument, so $S$ is
in the component of $V$. There are $\sum_{j\le w}\binom{|V|}{j}$ such sets.
For the second claim, $\att(V\setminus S)$ is the set of vertices outside
$S$ with a neighbour in $S$, which is $N(S)\setminus S$; so
$S\in\mathcal F_w$ if and only if $V\setminus S\in\mathcal F'_w$, and adding
a vertex to $S$ removes it from $V\setminus S$.
\end{proof}

The theorem places the free region, where every set is feasible before any
constraint bites, at the top of the lattice under the outer boundary and at
the bottom under the inner one, and says that a search avoiding it is the
primal search under the other name. What it does not say is how the two
components compare; that is a measurement, and it was made
(Section~\ref{sec:certs}, the two-ended enumeration): on 64 bound-gap
instances at 30 to 50 customers the component of $\emptyset$ has between
2,000 and 281,000 states, while the component of $V$ exceeded 3,000,000 on
61 of them and was 300 to 1,000 times the other on the three where it was
enumerated to the end. A dual search in this lattice would need a rule that
discards most of that component, of a kind no primal rule is: the rules of
Section~\ref{sec:search} bring the visited states to within a factor of
about two of the component, not orders of magnitude below it.

A second concrete dual search maximises over a compact class of frames. The
class we built, cliques, minor pullbacks and linked branches, has its ceiling
measured in Section~\ref{sec:certs}; here is the one statement about such
classes that is a theorem rather than a measurement.

\begin{proposition}[connected frames are blind to $\pw>\tw$]\label{prop:connected-frames}
Let $F_1,\dots,F_r$ be connected vertex sets with $|F_i|\ge n+1$, pairwise
inseparable by fewer than $n+1$ vertices, and such that no set of at most
$n$ vertices meets every $F_i$. Then the frame $F_1,\dots,F_r$ generates a
blockage of order $n$, and $\tw(G)\ge n$. In particular such a frame
certifies $\pw(G)\ge n$ only on graphs with $\tw(G)\ge n$.
\end{proposition}
\begin{proof}
For the blockage, check (a) and (b) of Lemma~\ref{lem:frame}. A separation
of order at most $n$ that splits every $F_j$ has, since each $F_j$ is
connected, a separator meeting every $F_j$, which is a set of at most $n$
vertices meeting all of them; so some $F_j$ lies inside a side, which is
(a). Two frame sets on opposite sides of a separation of order at most $n$
would be separated by at most $n$ vertices, against the hypothesis; this is
(b). For the treewidth, take a tree decomposition of width $\tw(G)$. Each
connected $F_i$ induces a subtree of the decomposition tree. If two of these
subtrees are disjoint, some bag separates $F_i$ from $F_j$, so
$\tw(G)+1\ge n+1$. Otherwise the subtrees pairwise intersect, and subtrees
of a tree have the Helly property, so some bag meets every $F_i$; that bag
has at most $\tw(G)+1$ vertices, so $\tw(G)+1\ge n+1$.
\end{proof}

This is the pathwidth copy of a bramble, and the proposition says it is only
a bramble: it certifies what treewidth certifies. The ingredient that sees
$\pw>\tw$ is branching, and the census of Section~\ref{sec:certs} measures
how rarely branching is the obstruction.

\subsection{What is open: the size of a blockage}

Theorem~\ref{thm:dual-conp} is conditional. For treewidth,
\citet{GroheMarx2009} proved the unconditional version in a fixed
representation: on bounded-degree expanders every bramble of order at least
$k^{1/2+\varepsilon}$ has $2^{\Omega(k^{\varepsilon})}$ sets, while every
graph of treewidth $k$ has a bramble of order
$\Omega(k^{1/2}/\log^2 k)$ with polynomially many sets; see also
\citet{HatzelKomosaPilipczukSorge2022}. The analogue for blockages, in the
frame representation of Lemma~\ref{lem:frame}, is the natural open
question, and we state it as a conjecture with the data beside it.

\begin{conjecture}[blockage size]\label{conj:blockage-size}
There is a family of bounded-degree graphs on $|V|$ vertices with
$\pw=\Theta(|V|)$ on which every blockage of order at least
$|V|^{1/2+\varepsilon}$ has a frame of $2^{\Omega(|V|^{\varepsilon})}$ sets;
and every graph has a blockage of order $\pw^{\Omega(1)}$ whose frame has
polynomially many sets.
\end{conjecture}

The second half follows, with a small exponent, from the certificate of
\citet{GroenlandJoretNadaraWalczak2023}, which is the iterated branch
operator above. For the first we have measurements and no proof. The
minimum frame size over all blockages of each order was computed exactly, by
a SAT model of the axioms, on every connected graph with at most nine
vertices and on structured and expander-like families to 22 vertices. At
order equal to the pathwidth the minimum is one on 28\% of the graphs on
nine vertices and at most three on 96\%, never above 13; on every graph with
$\pw>\tw$ it is at least three, on 13,548 of 13,548 such graphs; on random
cubic graphs it is flat in $|V|$ at fixed pathwidth and grows by a factor of
1.5 to 2 per unit of pathwidth, from 2 at pathwidth 3 to between 6 and 11 at
pathwidth 6; the largest values are 15 on the Petersen graph, 22 on the
Shrikhande graph and 32 on the Clebsch graph, about twice the number of
vertices. Two laws fit every exact number, a constant times $2^{\pw}$ and
$2|V|$, and with pathwidth at most nine they cannot be told apart. If the
first is the law, the conjecture's hard half follows on expanders, where
pathwidth is linear in $|V|$; the data show the factor per unit of order and
cannot see the exponent. A proof would need the expansion argument of
\citet{GroheMarx2009} redone for the invisible-fugitive game, since
strategies transfer only from the visible robber to the invisible fugitive
and not back. We found no prior statement about the size of a blockage.
